\documentclass[12pt,letterpaper]{article} % 12pt font, letter paper

% --- PACKAGES ---
\usepackage[utf8]{inputenc} % Input encoding
\usepackage[spanish]{babel} % Language support
\usepackage{times} % Use Times font (approximates Times New Roman)
\usepackage[margin=2.54cm]{geometry} % Margins 2.54 cm (1 inch)
\usepackage{setspace} % For line spacing
\usepackage{graphicx} % If you need images later
\usepackage{booktabs} % For professional tables
\usepackage{caption}  % For customizing captions
\usepackage{amsmath} % For math symbols if needed
\usepackage{amsfonts} % For math fonts if needed
\usepackage{amssymb} % For math symbols if needed
\usepackage{xurl} % Better URL breaking (replaces url package)
\usepackage[hidelinks]{hyperref} % For clickable URLs (optional)
\usepackage{ragged2e} % For RaggedRight justification (simulates flush left)
\usepackage{indentfirst} % Indent the first paragraph after a section heading
\usepackage{titlesec} % For customizing section titles
\usepackage{tocloft}
\usepackage{titletoc} % Added based on user template usage
\usepackage{enumitem} % For customizing lists (APA 7 indentation)

% --- APA 7 List Configuration (Hanging Indent + Circular Bullets) ---
% Note: Document uses \setstretch{2} globally, so lists inherit double spacing
% We set itemsep, parsep, topsep to 0pt to avoid EXTRA spacing beyond the global setting

% Itemize (Bulleted Lists)
\setlist[itemize]{
    leftmargin=1.9cm,      % Total left margin (1.27cm + 0.63cm space)
    itemindent=0pt,        % No additional indent for the bullet
    labelsep=0.63cm,       % Space between bullet and text
    label=\textbullet,     % Force circular bullets (not squares)
    itemsep=0pt,           % No extra spacing between items (uses global double spacing)
    parsep=0pt,            % No extra spacing between paragraphs within items
    topsep=0pt             % No extra spacing before/after list (uses global double spacing)
}

% Enumerate (Numbered Lists)
\setlist[enumerate]{
    leftmargin=1.9cm,      % Total left margin (1.27cm + 0.63cm space)
    itemindent=0pt,        % No additional indent for the number
    labelsep=0.63cm,       % Space between number and text
    itemsep=0pt,           % No extra spacing between items (uses global double spacing)
    parsep=0pt,            % No extra spacing between paragraphs within items
    topsep=0pt             % No extra spacing before/after list (uses global double spacing)
}

% --- Preamble Addition: Rename Table Label for Spanish ---
\addto\captionsspanish{\renewcommand{\tablename}{Tabla}}
\addto\captionsspanish{\renewcommand{\contentsname}{Tabla de Contenido}}
\addto\captionsspanish{\renewcommand{\listtablename}{Lista de Tablas}}

% --- APA 7 UNAD FORMATTING ---

\setstretch{2}

% Paragraph Indentation (0.5 inches = 1.27 cm)
\setlength{\parindent}{1.27cm}

% Disable Section Numbering (Fixes ToC Overlap)
\setcounter{secnumdepth}{0}

% Include Level 4 Headings (paragraph) in Table of Contents
\setcounter{tocdepth}{4}

% Page Numbering (Top Right)
\usepackage{fancyhdr}
\pagestyle{fancy}
\fancyhf{} % Clear header/footer
\renewcommand{\headrulewidth}{0pt} % No header rule
\fancyhead[R]{\thepage} % Page number top right
\fancypagestyle{plain}{% Redefine plain style used by titlepage etc.
  \fancyhf{}%
  \renewcommand{\headrulewidth}{0pt}%
  \fancyhead[R]{\thepage}%
}

% --- Preamble: Customization for Table of Contents (No Section Numbers in ToC) ---

% Remove numbering for sections in ToC
\renewcommand{\cftsecpresnum}{\relax} % No prefix for section number
\renewcommand{\cftsecnumwidth}{0pt}   % No space reserved for section number

% Remove numbering for subsections in ToC
\renewcommand{\cftsubsecpresnum}{\relax}
\renewcommand{\cftsubsecnumwidth}{0pt}

% Remove numbering for subsubsections in ToC
\renewcommand{\cftsubsubsecpresnum}{\relax}
\renewcommand{\cftsubsubsecnumwidth}{0pt}

% Ensure only the title and page number appear, with leaders (dots)
\renewcommand{\cftsecleader}{\cftdotfill{\cftdotsep}}
\renewcommand{\cftsubsecleader}{\cftdotfill{\cftdotsep}}
\renewcommand{\cftsubsubsecleader}{\cftdotfill{\cftdotsep}}

% Justification (Flush Left - Ragged Right)
\RaggedRight % Apply ragged right justification globally

% --- Customization for ToC and LoT Titles (Combined for Centering and Size) ---

% --- Para la Tabla de Contenido ---
\setlength{\cftbeforetoctitleskip}{0\baselineskip}
\setlength{\cftaftertoctitleskip}{0\baselineskip}

\renewcommand{\cfttoctitlefont}{\hfill\normalfont\bfseries\normalsize}
\renewcommand{\cftaftertoctitle}{\null\hfill}

% --- Para la Lista de Tablas ---
\setlength{\cftbeforelottitleskip}{0\baselineskip}
\setlength{\cftafterlottitleskip}{0\baselineskip}

\renewcommand{\cftlottitlefont}{\hfill\normalfont\bfseries\normalsize}
\renewcommand{\cftafterlottitle}{\null\hfill}

% --- Customization for ToC Entry Fonts (No Bold for Section Titles in ToC) ---
\renewcommand{\cftsecfont}{\normalfont}
\renewcommand{\cftsecpagefont}{\normalfont}
\renewcommand{\cftsubsecfont}{\normalfont}
\renewcommand{\cftsubsecpagefont}{\normalfont}
\renewcommand{\cftsubsubsecfont}{\normalfont}
\renewcommand{\cftsubsubsecpagefont}{\normalfont}

% --- Indentation for ToC entries ---
\setlength{\cftsecindent}{0em}
\setlength{\cftsubsecindent}{1.5em}
\setlength{\cftsubsubsecindent}{3em}

\newlength{\mytocskip}
\setlength{\mytocskip}{0\baselineskip}
\setlength{\cftbeforesecskip}{\mytocskip}
\setlength{\cftbeforesubsecskip}{\mytocskip}
\setlength{\cftbeforesubsubsecskip}{\mytocskip}

% Change the font for the table number in LoT to bold
\renewcommand{\cfttabpresnum}{\tablename~}
\renewcommand{\cfttabfont}{\bfseries}

% Configuration for List of Tables (LoT) entries
\titlecontents{table}
  [0em]
  {{\bfseries\tablename~\thecontentslabel\quad}\itshape}
  {}
  {}
  {\quad\cftdotfill{\cftdotsep}\normalfont\thecontentspage}
  [\vspace{0\baselineskip}]

% Section Title Formatting
\titleformat{\section}{\normalfont\bfseries\centering}{}{0em}{}
\titlespacing*{\section}{0pt}{0\baselineskip}{0\baselineskip}

\titleformat{\subsection}{\normalfont\bfseries}{}{0em}{}
\titlespacing*{\subsection}{0pt}{0\baselineskip}{0\baselineskip}

\titleformat{\subsubsection}{\normalfont\bfseries\itshape}{}{0em}{}
\titlespacing*{\subsubsection}{0pt}{0\baselineskip}{0\baselineskip}

% Level 4 Heading: Indented 1.27cm, Bold, Title Case, Period, Runin
\titleformat{\paragraph}[runin]{\normalfont\bfseries}{}{0em}{}[.]
\titlespacing*{\paragraph}{1.27cm}{0\baselineskip}{1em}

% Level 5 Heading: Indented 1.27cm, Bold Italic, Title Case, Period, Runin
\titleformat{\subparagraph}[runin]{\normalfont\bfseries\itshape}{}{0em}{}[.]
\titlespacing*{\subparagraph}{1.27cm}{0\baselineskip}{1em}

% Table Caption Formatting
\DeclareCaptionLabelFormat{table_num_bold_left}{\textbf{#1~#2}}
\DeclareCaptionLabelSeparator{mediumspace}{\par\vspace{0.5\baselineskip}}
\captionsetup[table]{
    labelformat=table_num_bold_left,
    labelsep=mediumspace,
    textfont=it,
    justification=raggedright,
    singlelinecheck=false,
    skip=\baselineskip,
    position=top
}

% --- DOCUMENT START ---
\begin{document}

% --- TITLE PAGE ---
\begin{titlepage}
    \thispagestyle{plain} % Force page number on title page
    \centering
    {\normalfont\bfseries Prototipo Web 3D Interactivo para la Visualización Técnica y Análisis Estructural de Hardware de Alto Rendimiento mediante Pipelines de Optimización Gráfica y WebAssembly\par}

    \vspace{2\baselineskip}
    \vspace{2\baselineskip}

    {\normalfont Estudiante\par}
    {\normalfont Alexander Woodcock Salomón\par}

    \vspace{2\baselineskip}
    \vspace{2\baselineskip}

    {\normalfont Asesor\par}
    {\normalfont Deivid Enrique Triviño Lozada\par}

    \vfill

    {\normalfont Universidad Nacional Abierta y a Distancia - UNAD\par
     Escuela de Ciencias Básicas, Tecnología e Ingeniería\par
     Ingeniería Multimedia\par}

    \vspace{1\baselineskip}
    {2025\par}
\end{titlepage}

% --- MAIN CONTENT ---
\pagestyle{fancy}

% --- RESUMEN ---
\section*{Resumen}
\addcontentsline{toc}{section}{Resumen}
\setlength{\parindent}{0pt} % No indent for Resumen
El diseño de hardware de alto rendimiento (drones, sistemas robóticos) enfrenta un desafío crítico en la comunicación técnica: los medios digitales tradicionales (imágenes 2D, PDFs estáticos) imponen alta carga cognitiva para comprender estructuras tridimensionales complejas. Este proyecto abordó esta brecha mediante el diseño, desarrollo y evaluación de un prototipo web 3D interactivo basado en Unity WebGL que optimiza la visualización técnica de hardware mediante \textit{pipelines} de optimización gráfica avanzados y WebAssembly. El marco teórico integró teoría de carga cognitiva (Mayer, Paivio), interacción humano-computador en 3D (Bowman, Norman), y renderizado físicamente basado (Cook-Torrance, Burley). La metodología siguió Design Science Research con desarrollo iterativo en Sprints de 4 semanas, implementando técnicas de retopología, \textit{baking} de normales, y Asset Bundles para cumplir con KPIs cuantitativos ($<$ 100,000 polígonos, $>$ 30 FPS en móvil mid-range, $<$ 3s shell / $<$ 10s carga completa). La validación incluyó \textit{profiling} de rendimiento (Unity Profiler), pruebas de usabilidad (SUS) y medición de carga cognitiva (NASA-TLX) con ingenieros/técnicos (N=8-12). El producto final es un MVP científico que valida la viabilidad técnica de WebGL para visualización de ingeniería, estableciendo un \textit{pipeline} replicable documentado según estándares académicos.

\hspace{1.27cm}\textit{\textbf{Palabras clave:}} WebGL, Unity, WebAssembly, Renderizado Físicamente Basado (PBR), Visualización 3D Interactiva, Optimización de Assets 3D, Carga Cognitiva, Interacción Humano-Computador, Technical Art, Hardware de Alto Rendimiento.

\newpage

% --- ABSTRACT ---
\section*{Abstract}
\addcontentsline{toc}{section}{Abstract}
\setlength{\parindent}{0pt} % No indent for Abstract
High-performance hardware design (drones, robotic systems) faces a critical challenge in technical communication: traditional digital media (2D images, static PDFs) impose high cognitive load for understanding complex three-dimensional structures. This project addressed this gap through the design, development, and evaluation of an interactive 3D web prototype based on Unity WebGL that optimizes technical hardware visualization through advanced graphics optimization \textit{pipelines} and WebAssembly. The theoretical framework integrated cognitive load theory (Mayer, Paivio), 3D human-computer interaction (Bowman, Norman), and physically-based rendering (Cook-Torrance, Burley). The methodology followed Design Science Research with iterative development in 4-week Sprints, implementing retopology, normal \textit{baking}, and Asset Bundles to meet quantitative KPIs ($<$ 100,000 polygons, $>$ 30 FPS on mid-range mobile, $<$ 3s shell / $<$ 10s full load). Validation included performance \textit{profiling} (Unity Profiler), usability testing (SUS), and cognitive load measurement (NASA-TLX) with engineers/technicians (N=8-12). The final product is a scientific MVP that validates the technical feasibility of WebGL for engineering visualization, establishing a replicable, academically documented \textit{pipeline}.

\hspace{1.27cm}\textit{\textbf{Keywords:}} WebGL, Unity, WebAssembly, Physically-Based Rendering (PBR), Interactive 3D Visualization, 3D Asset Optimization, Cognitive Load, Human-Computer Interaction, Technical Art, High-Performance Hardware.

\newpage

% Restore indentation for the rest of the document
\setlength{\parindent}{1.27cm}

% --- TABLE OF CONTENTS ---
\newpage
\tableofcontents
\newpage

% --- LIST OF TABLES ---
\newpage
\renewcommand{\listtablename}{Lista de Tablas}
\listoftables
\newpage

% --- INFORMACIÓN GENERAL ---
\newpage
\section*{Información General del Trabajo de Grado}
\addcontentsline{toc}{section}{Información General del Trabajo de Grado}

\subsection*{Integrantes del Proyecto}
\addcontentsline{toc}{subsection}{Integrantes del Proyecto}

\textit{Nombre del estudiante:} Alexander Woodcock Salomón

\textit{Identificación C.C.:} 1018474351

\textit{Programa Académico:} Ingeniería Multimedia

\textit{No. de Créditos Aprobados:} 152

\textit{\% de créditos aprobados:} 92.11

\textit{Correo electrónico:} awoodcocks@unadvirtual.edu.co

\textit{Teléfono / Celular:} 3054396581

\textit{Dirección residencia:} Finca Armenia casa 2

\textit{Municipio / Departamento:} Pasto / Nariño

\textit{CENTRO:} ZCSUR

\textit{ZONA:} PASTO

\subsection*{Datos Específicos del Proyecto}
\addcontentsline{toc}{subsection}{Datos Específicos del Proyecto}

\textit{Duración del proyecto (meses):} Seis (6)

\textit{Línea de Investigación:} Ingeniería Multimedia, Visualización Técnica 3D Web

\textit{Escuela:} Escuela de Ciencias Básicas, Tecnología e Ingeniería (ECBTI)

\textit{Descriptores palabras clave:} WebGL, Unity, WebAssembly, PBR, Visualización 3D, Optimización, Carga Cognitiva, HCI, Technical Art.

\newpage

% --- SECTIONS (Reutilizando las de la propuesta pero adaptadas mentalmente a pasado, 
%     se recomienda crear copias en Informe_final/sections si se van a editar masivamente) ---
%     Para este paso inicial, usaremos las secciones existentes pero añadiremos las nuevas.

\section{Planteamiento del Problema}

En la documentaci\'on de hardware t\'ecnico complejo, la comunicaci\'on visual sigue anclada con frecuencia en paradigmas est\'aticos, como archivos PDF, planos 2D y l\'aminas de ensamblaje, que disocian la informaci\'on visual de su contexto espacial. En esta propuesta, \textit{hardware complejo} se refiere a sistemas electromec\'anicos multicomponente con integraci\'on estructural, energ\'etica, electr\'onica y de control, alta densidad de relaciones espaciales entre piezas y necesidad de inspecci\'on visual detallada del ensamblaje. Esta desconexi\'on obliga a ingenieros y t\'ecnicos a realizar procesos mentales de reconstrucci\'on tridimensional que pueden incrementar la demanda cognitiva asociada a la tarea y consumir recursos de memoria de trabajo (Sweller, 1988).

Sin embargo, la migraci\'on hacia entornos 3D interactivos en la web enfrenta restricciones arquitect\'onicas significativas. Los modelos CAD de ingenier\'ia, basados en NURBS, ensamblajes extensos y activos de alta densidad geom\'etrica, no son directamente compatibles con las condiciones de despliegue de una aplicaci\'on Web orientada a acceso desde navegador. Aunque WebAssembly admite crecimiento din\'amico de memoria, la viabilidad real de Unity Web depende de controlar la huella inicial de los activos, la fragmentaci\'on del espacio de memoria del navegador y la necesidad de reservar bloques contiguos suficientemente grandes para heap y datos de aplicaci\'on (Unity Technologies, 2024a; WebAssembly Community Group, s. f.). En este contexto, persiste la necesidad de documentar y evaluar un \textit{pipeline} de \textit{Technical Art} que permita cerrar la brecha entre fidelidad visual, claridad t\'ecnica y rendimiento web sin sacrificar la legibilidad funcional del ensamblaje.

El problema t\'ecnico puede resumirse en la relaci\'on entre fidelidad y rendimiento:
\begin{itemize}
    \item \textbf{Huella inicial de datos.} Modelos CAD originales pueden alcanzar decenas o cientos de megabytes, y crecer a escalas mayores en ensamblajes extensos, por lo que requieren limpieza, simplificaci\'on geom\'etrica, \textit{baking} y empaquetado selectivo para ser viables en web.
    \item \textbf{Rendimiento en tiempo real.} Mantener 30 FPS o m\'as, equivalente a un \textit{frame time} de aproximadamente 33.33 ms, exige controlar presupuesto geom\'etrico, materiales, \textit{draw calls} y estados de render.
    \item \textbf{Latencia de interacci\'on.} Soluciones basadas en \textit{streaming} remoto pueden ofrecer alta fidelidad, pero introducen dependencia de red y latencias poco adecuadas para inspecci\'on t\'ecnica precisa.
\end{itemize}

Existe, por tanto, la necesidad de dise\~nar y evaluar un \textit{pipeline} de \textit{Technical Art} que permita desplegar modelos de hardware complejo en un visor 3D web con orientaci\'on espacial suficiente para tareas de inspecci\'on, costo computacional controlado y capacidad de an\'alisis visual del ensamblaje para prop\'ositos acad\'emicos e ingenieriles.

\subsection{Pregunta de investigaci\'on}

\noindent \textbf{Pregunta principal.} \textquestiondown Qu\'e diferencias descriptivas se observan en desempe\~no de tareas, usabilidad percibida y carga de trabajo percibida cuando usuarios con perfil af\'in realizan actividades de exploraci\'on e identificaci\'on del dron Holybro X500~V2 mediante un prototipo Web 3D interactivo, en comparaci\'on con un soporte 2D de referencia, y bajo qu\'e condiciones de rendimiento resulta t\'ecnicamente viable dicho prototipo en el navegador?

\newpage

\section{Justificaci\'on}

La elecci\'on de Unity WebGL sobre librer\'ias nativas de JavaScript (Three.js, Babylon.js) o soluciones de \textit{streaming} remoto (Unreal Pixel Streaming) responde a requisitos concretos de visualizaci\'on t\'ecnica, desempe\~no web y viabilidad del \textit{pipeline} de arte t\'ecnico. El objetivo del proyecto no es solo mostrar un modelo 3D, sino construir un prototipo navegable, mantenible y evaluable acad\'emicamente, con interacciones t\'ecnicas, modos anal\'iticos y documentaci\'on consistente.

\subsection{Criterios t\'ecnicos de selecci\'on}

\subsubsection{Ejecuci\'on de l\'ogica en WebAssembly}

Unity utiliza IL2CPP para transpilar c\'odigo C\# a C++ y posteriormente a WebAssembly (WASM). Esto permite ejecutar la l\'ogica de interacci\'on y orquestaci\'on del prototipo en un entorno compilado y predecible dentro del navegador, con una arquitectura m\'as robusta para sistemas con m\'ultiples managers, estados de interfaz y herramientas de visualizaci\'on t\'ecnica.

\subsubsection{\textit{Pipeline} integrado de importaci\'on y renderizado}

A diferencia de un enfoque puramente manual en bibliotecas de bajo nivel, Unity ofrece un entorno integrado para importar geometr\'ia, gestionar materiales, perfilar rendimiento y desplegar un build WebGL desde una misma base de trabajo. Esta integraci\'on es especialmente valiosa cuando el proyecto requiere coordinar limpieza de activos CAD, shaders personalizados, UI responsiva y pruebas funcionales sin fragmentar el flujo entre demasiadas herramientas.

\subsubsection{Capacidad de extensi\'on para visualizaci\'on t\'ecnica}

El proyecto exige combinar renderizado realista con modos de inspecci\'on y an\'alisis. Unity URP permite aprovechar un flujo PBR consolidado y, al mismo tiempo, extender el sistema con shaders personalizados para corte transversal, vista \textit{X-Ray}, visualizaci\'on t\'ermica y otros modos anal\'iticos. Esta flexibilidad es determinante porque la versi\'on final del prototipo no depende de una sola visualizaci\'on, sino de una familia de modos con prop\'ositos distintos.

\subsubsection{Herramientas para colaboraci\'on entre programaci\'on y arte t\'ecnico}

El flujo de trabajo requiere colaboraci\'on entre programadores y artistas t\'ecnicos. Unity facilita esta articulaci\'on mediante materiales editables, inspecci\'on visual del contenido, herramientas de editor y un ecosistema orientado a iteraci\'on r\'apida. Esto reduce fricci\'on en tareas como ajuste de mallas, configuraci\'on de materiales, depuraci\'on de escenas y afinaci\'on de la experiencia visual.

\subsubsection{Escalabilidad y pertinencia disciplinar}

El uso de un motor ampliamente adoptado permite que el prototipo sea mantenible, replicable y escalable. Desde la perspectiva de Ingenier\'ia Multimedia, esta decisi\'on tambi\'en es pertinente porque conecta modelado 3D, optimizaci\'on gr\'afica, dise\~no de interfaz, evaluaci\'on con usuarios y publicaci\'on web en un mismo caso de estudio.

\subsubsection{Ventana de oportunidad tecnol\'ogica}

La convergencia entre WebAssembly, WebGL 2.0 y \textit{pipelines} modernos de renderizado abre una oportunidad concreta para democratizar el acceso a visualizaciones t\'ecnicas interactivas desde el navegador. El valor del proyecto no radica solo en optimizar recursos, sino en demostrar que una aplicaci\'on web puede soportar un caso de visualizaci\'on t\'ecnica con criterios reales de claridad espacial, usabilidad y rendimiento.

\subsection{Aporte a la Ingenier\'ia Multimedia}
\begin{itemize}
    \item documentar un \textit{pipeline} replicable para llevar activos CAD a un visor WebGL interactivo;
    \item validar el rol del ingeniero multimedia como puente entre complejidad t\'ecnica y experiencia de usuario;
    \item generar criterios verificables de rendimiento y consistencia funcional para visualizaci\'on web 3D;
    \item articular fundamentos de computaci\'on gr\'afica, UX y carga cognitiva en un proyecto aplicado.
\end{itemize}

\newpage

\section{Objetivos}

\subsection{Objetivo General}
Desarrollar un prototipo de visualizaci\'on Web 3D interactiva basado en Unity WebGL, enfocado en la exploraci\'on t\'ecnica y el an\'alisis visual de hardware de alto rendimiento, optimizado mediante criterios de \textit{Technical Art} para su despliegue eficiente en navegadores web est\'andar.

\subsection{Objetivos Espec\'ificos}

Dise\~nar un \textit{pipeline} de optimizaci\'on de activos 3D que reduzca la carga poligonal de modelos CAD originales hasta un presupuesto geom\'etrico compatible con WebGL, manteniendo fidelidad visual mediante retopolog\'ia, limpieza de mallas y \textit{baking} de texturas.

Integrar en Unity URP un sistema de materiales y modos de visualizaci\'on que combine renderizado realista con herramientas de inspecci\'on t\'ecnica, procurando un \textit{frame time} igual o inferior a 33.33 ms (30 FPS o m\'as) en el dispositivo objetivo.

Programar un sistema de interacci\'on en C\# compilado a WebAssembly que permita navegaci\'on orbital, selecci\'on de piezas, consulta contextual de informaci\'on y uso de herramientas como vista explosionada, corte transversal y modos anal\'iticos del prototipo final.

Validar el rendimiento, la usabilidad y la carga cognitiva del prototipo mediante herramientas de perfilado, pruebas con usuarios, cuestionarios SUS, NASA-TLX Raw y protocolo \textit{Think-Aloud}, contrastando los resultados frente a las metas operativas definidas para el proyecto.

\newpage

\section{Marco de Referencia}
\subsection{Estado del Arte}

El estado del arte analiza tecnolog\'ias actuales de visualizaci\'on 3D en la web, evaluando sus capacidades, limitaciones y pertinencia para la visualizaci\'on t\'ecnica de hardware complejo. Se examinan bibliotecas nativas, motores de juego, plataformas \textit{SaaS} y soluciones de \textit{streaming} para identificar la herramienta m\'as adecuada en t\'erminos de control del \textit{pipeline}, capacidad de extensi\'on y viabilidad de desarrollo.

\subsubsection{Benchmarking de soluciones Web 3D}

La visualizaci\'on 3D en la web ha evolucionado significativamente desde la introducci\'on de WebGL en 2011 (Khronos Group, 2011). Yu et al. (2023) identifican WebGL como la tecnolog\'ia dominante para renderizado 3D en navegadores, con creciente competencia de WebGPU. A continuaci\'on se presentan las principales alternativas evaluadas:

\paragraph{Three.js.} Biblioteca ligera construida sobre WebGL, ampliamente adoptada por la comunidad. Destaca por su curva de aprendizaje accesible, tama\~no de \textit{bundle} relativamente reducido y documentaci\'on extensa. Como desventaja, exige construir de forma manual gran parte del \textit{pipeline} de importaci\'on, gesti\'on de memoria e integraci\'on avanzada de herramientas.

\paragraph{Babylon.js.} Motor WebGL/WebGPU con enfoque \textit{enterprise}, soporte para GUI, f\'isica y escenas complejas. Ofrece una base muy completa, pero el esfuerzo de configuraci\'on para un flujo t\'ecnico altamente personalizado sigue siendo considerable.

\paragraph{Unity WebGL.} Motor con compilaci\'on a Web mediante IL2CPP, WebAssembly y el \textit{toolchain} del build. Su ecosistema integra editor visual, profiler, \textit{pipeline} de materiales, herramientas para UI y soporte para extensi\'on mediante c\'odigo C\# y shaders personalizados (Unity Technologies, 2024b, 2024d). La documentaci\'on actual de Unity 6 contempla soporte para navegadores compatibles de escritorio y algunos navegadores m\'oviles, aunque el comportamiento efectivo depende del navegador, la versi\'on, la memoria disponible y el nivel de optimizaci\'on del proyecto (Unity Technologies, s. f.). Su principal costo es una huella inicial de build mayor que la de bibliotecas puramente web.

\paragraph{Unreal Engine Pixel Streaming.} Soluci\'on de \textit{streaming} remoto que ofrece muy alta fidelidad visual a cambio de depender de infraestructura GPU del lado del servidor. Aunque es potente para demostraciones fotorrealistas, su latencia y costo operativo la hacen menos adecuada para un prototipo acad\'emico de acceso web directo.

\paragraph{Spline.} Herramienta \textit{web-first} orientada a dise\~no iterativo y colaboraci\'on. Es apropiada para presentaci\'on de piezas 3D ligeras y experiencias visuales sencillas, pero no est\'a pensada para sistemas t\'ecnicos con l\'ogica de interacci\'on compleja.

\paragraph{Marmoset Viewer.} Visor especializado en presentaci\'on PBR de alta calidad. Su fortaleza es la exhibici\'on visual del modelo, pero su margen de personalizaci\'on funcional es limitado para casos de inspecci\'on t\'ecnica.

\paragraph{Sketchfab.} Plataforma de \textit{hosting} y visualizaci\'on 3D con soporte de anotaciones y publicaci\'on r\'apida. Aunque es conveniente para difusi\'on, presenta menor control sobre arquitectura, l\'ogica e integraci\'on profunda de herramientas personalizadas.

\paragraph{PlayCanvas.} Plataforma de desarrollo web con editor colaborativo y \textit{runtime} ligero. Resulta atractiva para proyectos nativos del navegador, pero su ecosistema y documentaci\'on para casos t\'ecnicos muy espec\'ificos siguen siendo menores que los de Unity.

\begin{table}[H]
\centering
\caption{Comparativa cualitativa de soluciones Web 3D}
\small
\resizebox{\textwidth}{!}{%
\begin{tabular}{lcccccc}
\hline
\textbf{Tecnolog\'ia} & \textbf{Huella inicial} & \textbf{Control del pipeline} & \textbf{PBR} & \textbf{Interactividad} & \textbf{Infraestructura externa} & \textbf{Pertinencia} \\ \hline
Three.js & Baja & Alto, pero manual & Manual/extendible & Total & No & Alta para desarrollo web a medida \\
Babylon.js & Media & Alto & Integrado & Total & No & Alta para apps web 3D complejas \\
Unity WebGL & Alta & Alto e integrado & URP & Total & No & Alta para prototipos t\'ecnicos con tooling \\
UE Pixel Streaming & Media/Alta & Alto & Muy alto & Alta con latencia & S\'i & Media para acceso web directo \\
Spline & Baja & Bajo & B\'asico & Limitada & No & Media para presentaci\'on ligera \\
Marmoset Viewer & Baja/Media & Bajo & Alto & Baja/Media & No & Media para exhibici\'on visual \\
Sketchfab & Media & Medio & Bueno & Media & Parcial & Media para publicaci\'on r\'apida \\
PlayCanvas & Baja & Alto & Integrado & Total & No & Alta para experiencias web nativas \\ \hline
\end{tabular}
}
\par\vspace{0.5\baselineskip}
\textit{Nota}. La tabla presenta una comparaci\'on cualitativa construida a partir de documentaci\'on oficial y literatura revisada; no corresponde a un benchmark uniforme ejecutado experimentalmente por el autor bajo un mismo protocolo. \textit{Fuente}. Elaboraci\'on propia con base en documentaci\'on oficial y literatura revisada.
\end{table}

\subsubsection{Brecha identificada y selecci\'on de Unity WebGL}

La evaluaci\'on t\'ecnica identific\'o limitaciones cr\'iticas en tres grupos de soluciones:

\begin{itemize}
    \item \textbf{Pixel Streaming.} Latencia de red y dependencia de infraestructura \textit{cloud} poco compatibles con un prototipo acad\'emico de acceso directo en navegador.
    \item \textbf{Spline, Marmoset y Sketchfab.} Restricciones para integrar l\'ogica de inspecci\'on t\'ecnica, flujo de detalle, modos de an\'alisis y arquitectura propia del sistema.
    \item \textbf{Three.js, Babylon.js y PlayCanvas.} Mayor esfuerzo de implementaci\'on para construir desde cero un \textit{pipeline} integral de importaci\'on, materiales, controles, clipping global, modos de visualizaci\'on t\'ecnica y herramientas auxiliares de auditor\'ia.
\end{itemize}

Unity WebGL fue seleccionado no por ofrecer la huella inicial m\'as baja, sino por el equilibrio entre cuatro factores:
\begin{itemize}
    \item \textbf{Integraci\'on del tooling.} Permite concentrar importaci\'on, materiales, UI, perfilado y build en un mismo entorno.
    \item \textbf{Herramientas visuales y de editor.} Facilita el trabajo conjunto entre programaci\'on, dise\~no y arte t\'ecnico.
    \item \textbf{Arquitectura tipada y extensible.} La cadena C\# $\rightarrow$ IL2CPP $\rightarrow$ WASM ofrece una base estructurada para la l\'ogica interactiva del prototipo (Unity Technologies, 2024d).
    \item \textbf{Soporte para modos anal\'iticos.} Reduce el riesgo t\'ecnico al implementar corte transversal, visualizaci\'on t\'ermica, vista \textit{X-Ray} y otras herramientas especializadas dentro del tiempo acad\'emico disponible.
\end{itemize}

\newpage

\subsection{Marco Te\'orico}

El dise\~no, desarrollo y evaluaci\'on de la plataforma web 3D se fundamenta en un marco te\'orico multidisciplinario que integra conocimientos de ingenier\'ia, dise\~no, psicolog\'ia cognitiva e inform\'atica gr\'afica. El proyecto no se limita a mostrar un modelo tridimensional: busca demostrar que una visualizaci\'on t\'ecnica interactiva puede mejorar la comprensi\'on espacial del hardware, sostener un rendimiento web aceptable y articular una experiencia de usuario m\'as clara que la documentaci\'on 2D est\'atica.

\paragraph{Hip\'otesis de usabilidad y carga cognitiva.} La visualizaci\'on 3D interactiva gestiona la carga cognitiva intr\'inseca mediante segmentaci\'on del contenido y reduce la carga extr\'inseca al externalizar operaciones de rotaci\'on mental que en soportes 2D deben resolverse internamente por el usuario. Esta hip\'otesis se validar\'a mediante tareas representativas, KPIs t\'ecnicos, SUS, NASA-TLX y observaci\'on \textit{Think-Aloud}.

El uso de una estrategia tipo MVP justifica la exclusi\'on intencional de caracter\'isticas no esenciales para la validaci\'on acad\'emica inmediata, como persistencia avanzada, instrumentaci\'on anal\'itica compleja o variantes de producto no integradas en la versi\'on final del prototipo.

\subsubsection{Teor\'ia de la Carga Cognitiva (Sweller, 1988; Sweller et al., 2019)}

La Teor\'ia de la Carga Cognitiva plantea que la memoria de trabajo tiene capacidad limitada y que el dise\~no de materiales e interfaces debe regular la demanda que impone sobre ella. Seg\'un Cowan (2001), esta capacidad operativa ronda unos pocos elementos simult\'aneos, por lo que la forma de presentar informaci\'on t\'ecnica incide directamente en la comprensi\'on y en la fatiga del usuario.

Sweller, van Merri\"enboer y Paas (2019) distinguen tres tipos de carga:

\paragraph{Carga intr\'inseca.} Corresponde a la complejidad inherente del contenido. En este proyecto es elevada porque el dron integra componentes estructurales, electr\'onicos y mec\'anicos con relaciones espaciales y funcionales entre s\'i.

\paragraph{Carga extr\'inseca.} Es la carga innecesaria generada por una presentaci\'on deficiente. En documentaci\'on 2D est\'atica, el usuario debe imaginar relaciones tridimensionales no visibles, lo cual incrementa la necesidad de rotaci\'on mental y consume recursos cognitivos que no est\'an dirigidos al objetivo t\'ecnico principal.

\paragraph{Carga germana.} Es el esfuerzo que s\'i contribuye a construir esquemas mentales \'utiles. En una interfaz 3D bien dise\~nada, los recursos cognitivos liberados por la reducci\'on de carga extr\'inseca pueden redirigirse a comprender mejor la funci\'on, ubicaci\'on y relaci\'on entre componentes.

En el contexto del prototipo, la interfaz 3D orbital y la consulta contextual de piezas buscan precisamente reducir la carga extr\'inseca y favorecer una carga germana orientada a comprensi\'on estructural.

\subsubsection{Navegaci\'on espacial e interacci\'on 3D}

Bowman et al. (2004) definen la interacci\'on 3D como aquella en la cual las tareas del usuario se ejecutan directamente en un contexto espacial tridimensional. Esto es relevante porque el valor del prototipo depende de que el usuario pueda explorar el ensamblaje sin perder contexto global mientras inspecciona detalles locales.

Los controles orbitales siguen principios de navegaci\'on espacial orientados a prevenir desorientaci\'on:
\begin{itemize}
    \item rotaci\'on alrededor de un foco de inter\'es estable;
    \item \textit{pan} y \textit{zoom} como operaciones separadas de la rotaci\'on;
    \item recentrado del inter\'es cuando se selecciona una pieza o se invoca una herramienta de inspecci\'on;
    \item continuidad de movimiento suficiente para evitar saltos bruscos de interpretaci\'on espacial.
\end{itemize}

\paragraph{Cognici\'on distribuida (Hutchins, 1995).} La interfaz no solo presenta informaci\'on: tambi\'en funciona como artefacto cognitivo. La c\'amara orbital externaliza parte de la rotaci\'on mental, los \textit{hotspots} y fichas de pieza operan como memoria externa y los modos anal\'iticos convierten procesos de inspecci\'on complejos en operaciones visuales guiadas.

\subsubsection{Heur\'isticas de usabilidad (Nielsen, 1994)}

Las heur\'isticas de Nielsen orientan la construcci\'on de la interfaz:
\begin{enumerate}
    \item \textbf{Visibilidad del estado del sistema.} La aplicaci\'on debe comunicar con claridad qu\'e pieza est\'a seleccionada, qu\'e modo visual est\'a activo y cu\'ando una herramienta de an\'alisis modifica la escena.
    \item \textbf{Correspondencia con el mundo real.} Los nombres de componentes, la organizaci\'on por categor\'ias y la l\'ogica de inspecci\'on deben alinearse con la terminolog\'ia t\'ecnica del dron.
    \item \textbf{Control y libertad del usuario.} La navegaci\'on, la salida de paneles y el reinicio de estado deben ser previsibles y reversibles.
    \item \textbf{Reconocimiento antes que memoria.} La interfaz debe reducir la necesidad de recordar nombres, ubicaciones o secuencias de interacci\'on, apoy\'andose en ayudas visuales y estados expl\'icitos.
\end{enumerate}

\subsubsection{Rotaci\'on mental y visualizaci\'on espacial}

Hegarty y Waller (2004) muestran que la interpretaci\'on de estructuras tridimensionales depende de habilidades espaciales que no se distribuyen homog\'eneamente entre usuarios. Bartlett y Dorribo Camba (2023) refuerzan que los modelos 3D interactivos pueden disminuir esta barrera al facilitar interpretaci\'on gr\'afica y comprensi\'on estructural. En consecuencia, el visor no se justifica solo por atractivo visual, sino como apoyo cognitivo para usuarios con distintos niveles de habilidad espacial.

\subsubsection{Percepci\'on visual y organizaci\'on de la informaci\'on}

Los principios de Gestalt (Wertheimer, 1923; K\"ohler, 1929) ayudan a estructurar la escena y la interfaz:
\begin{itemize}
    \item \textbf{Proximidad.} Elementos cercanos se perciben como pertenecientes al mismo subsistema.
    \item \textbf{Similitud.} El color o tratamiento visual de \textit{hotspots} y categor\'ias facilita asociaci\'on funcional.
    \item \textbf{Figura-fondo.} El modelo principal debe destacar sobre el fondo para no competir con el contenido t\'ecnico.
\end{itemize}

Ware (2012) y Miller (1956) tambi\'en sustentan el uso de categor\'ias crom\'aticas distinguibles, evitando esquemas que dificulten accesibilidad o saturen la memoria inmediata del usuario.

\subsubsection{Fundamentos matem\'aticos de optimizaci\'on gr\'afica}

\paragraph{Presupuesto poligonal.}
Una primera restricci\'on de complejidad geom\'etrica puede expresarse como:
\[
P_{\text{total}} = \sum_{i=1}^{n} P_i
\]
donde $P_i$ representa los tri\'angulos aportados por cada componente del modelo. En t\'erminos de planeaci\'on, el proyecto fija un presupuesto objetivo para mantener la escena dentro de un rango compatible con WebGL. Este presupuesto no es una ley universal del motor, sino un criterio de control para el proceso de optimizaci\'on.

\paragraph{Densidad de texel.}
La densidad de texel busca que la nitidez percibida de las texturas sea coherente entre componentes. Una forma operativa de expresarla es:
\[
\rho = \frac{R_{\text{texture}}}{L_{\text{mundo}}}
\]
donde $R_{\text{texture}}$ es la resoluci\'on efectiva aplicada sobre una direcci\'on o tramo de referencia, y $L_{\text{mundo}}$ es la longitud real correspondiente en el modelo. Si $\rho$ cambia bruscamente entre piezas, algunas se perciben borrosas y otras innecesariamente costosas en memoria.

\paragraph{Costo de dibujo y agrupaci\'on de mallas.}
El costo de renderizado no depende solo del n\'umero de tri\'angulos, sino tambi\'en del n\'umero de lotes o \textit{draw calls}. Sin agrupaci\'on, una aproximaci\'on simple es:
\[
D_{\text{base}} \approx \sum_{k=1}^{m} n_k
\]
donde $n_k$ representa la cantidad de objetos que requieren una combinaci\'on distinta de malla, material o pasada. Si existen condiciones de \textit{batching} o \textit{instancing} compatibles, varias instancias pueden reducirse a menos lotes efectivos:
\[
D_{\text{efectivo}} \approx \sum_{k=1}^{m} b_k, \qquad 1 \leq b_k \leq n_k
\]
Por tanto, la reducci\'on de \textit{draw calls} no debe modelarse como una divisi\'on exacta universal, sino como una funci\'on de compatibilidad entre materiales, mallas, pasadas de render y estado del \textit{pipeline}.

\paragraph{\textit{Frame time} y objetivo de FPS.}
En una aplicaci\'on interactiva, el tiempo por cuadro determina la fluidez percibida. Una aproximaci\'on \'util para perfilado es:
\[
T_{\text{frame}} \approx \max(T_{\text{CPU}}, T_{\text{GPU}})
\]
con un posible costo adicional por sincronizaci\'on. Para sostener 30 FPS, se requiere aproximadamente:
\[
T_{\text{frame}} \leq 33.33 \text{ ms}
\]
Esto traduce el rendimiento a una restricci\'on directamente medible en navegador y en herramientas de perfilado.

\subsubsection{Modelo matem\'atico de renderizado f\'isicamente basado (PBR)}

El renderizado f\'isicamente basado parte de la ecuaci\'on de render:
\[
L_o(\omega_o) = \int_{\Omega} f_r(\omega_i,\omega_o)\, L_i(\omega_i)\, (\mathbf{n}\cdot\omega_i)\, d\omega_i
\]
donde $L_o$ es la radiancia saliente en la direcci\'on de observaci\'on $\omega_o$, $L_i$ es la radiancia incidente desde la direcci\'on $\omega_i$, $\mathbf{n}$ es la normal de la superficie y $f_r$ es la BRDF del material. Esta ecuaci\'on expresa que el color visible depende de cu\'anta luz llega a la superficie, c\'omo la superficie la redistribuye y qu\'e parte de esa redistribuci\'on coincide con la direcci\'on de la c\'amara.

\paragraph{BRDF especular de Cook-Torrance.}
La componente especular puede modelarse como:
\[
f_{\text{spec}} = \frac{D(\mathbf{n},\mathbf{h})\, F(\mathbf{v},\mathbf{h})\, G(\mathbf{n},\mathbf{l},\mathbf{v})}{4(\mathbf{n}\cdot\mathbf{l})(\mathbf{n}\cdot\mathbf{v})}
\]
donde $\mathbf{l}$ es la direcci\'on de la luz, $\mathbf{v}$ la direcci\'on de vista y $\mathbf{h}$ el vector mitad entre ambas. Cada t\'ermino cumple una funci\'on distinta:
\begin{itemize}
    \item $D$ describe la distribuci\'on estad\'istica de microfacetas orientadas para reflejar luz hacia la c\'amara;
    \item $F$ modela el efecto Fresnel, es decir, el incremento de reflectancia cuando la observaci\'on se acerca al \'angulo rasante;
    \item $G$ aten\'ua la respuesta especular por auto-sombreado y enmascaramiento geom\'etrico entre microfacetas.
\end{itemize}

\paragraph{Aproximaci\'on de Fresnel de Schlick.}
Para reducir costo computacional, el t\'ermino de Fresnel suele aproximarse mediante:
\[
F = F_0 + (1 - F_0)(1 - \cos\theta)^5
\]
donde $F_0$ es la reflectancia a incidencia normal y $\theta$ es el \'angulo entre la direcci\'on de vista y el vector mitad. Esta aproximaci\'on permite obtener un comportamiento visualmente cre\'ible con un costo bajo en tiempo real.

\paragraph{Distribuci\'on GGX.}
La distribuci\'on de microfacetas adoptada com\'unmente en PBR en tiempo real es GGX:
\[
D_{\text{GGX}} = \frac{\alpha^2}{\pi\left((\mathbf{n}\cdot\mathbf{h})^2(\alpha^2 - 1) + 1\right)^2}
\]
donde $\alpha$ controla la amplitud del l\'obulo especular y suele derivarse de la rugosidad como $\alpha = roughness^2$. Valores bajos producen reflejos concentrados y valores altos reflejos anchos y difusos.

\paragraph{Separaci\'on difusa-especular y conservaci\'on de energ\'ia.}
En un flujo \textit{metallic-roughness}, la energ\'ia reflejada debe mantenerse acotada. De manera simplificada:
\[
k_d + k_s \leq 1
\]
donde $k_d$ representa la fracci\'on difusa y $k_s$ la especular. En materiales met\'alicos, la componente difusa tiende a desaparecer y el color especular toma protagonismo; en diel\'ectricos ocurre lo contrario.

\paragraph{Relaci\'on con la app real.}
En la aplicaci\'on final, el modo \texttt{Realistic} no reimplementa manualmente toda la BRDF microfacetada desde cero. El proyecto delega el c\'alculo PBR base al \textit{pipeline} de Unity URP mediante la funci\'on \texttt{UniversalFragmentPBR}, mientras que los shaders propios a\~naden l\'ogica espec\'ifica del sistema, como planos de corte globales y modos de inspecci\'on. Por ello, las ecuaciones anteriores se presentan como fundamento matem\'atico del modelo de iluminaci\'on empleado por URP y no como afirmaci\'on de un solver BRDF completamente escrito desde cero dentro del proyecto.

\paragraph{Iluminaci\'on basada en im\'agenes (IBL).}
La iluminaci\'on ambiental en PBR se apoya en mapas de entorno precalculados. Conceptualmente, la integral ambiental se aproxima mediante dos componentes: un mapa de irradiancia para la parte difusa y un mapa de entorno prefiltrado para la parte especular. Esto permite aproximar reflejos y contribuci\'on ambiental con costo constante por p\'ixel, sin necesidad de \textit{ray tracing} completo en tiempo real.

\subsubsection{Matem\'atica aplicada a shaders de inspecci\'on y an\'alisis}

\paragraph{Plano de corte global.}
El corte transversal de la app puede expresarse con la ecuaci\'on de un plano:
\[
\Pi = (n_x, n_y, n_z, d), \qquad d = -\mathbf{n}\cdot\mathbf{p}
\]
donde $\mathbf{n}$ es la normal del plano y $\mathbf{p}$ un punto perteneciente a dicho plano. Para cualquier punto del espacio $\mathbf{x}$, la distancia algebraica al plano se eval\'ua como:
\[
\mathbf{n}\cdot\mathbf{x} + d
\]
Si este valor es negativo, el fragmento queda del lado descartado del plano y no se renderiza:
\[
\mathbf{n}\cdot\mathbf{x} + d < 0 \Rightarrow \text{fragmento descartado}
\]
Esta formulaci\'on coincide con el mecanismo usado por los shaders del sistema cuando reciben los planos globales calculados por el gestor de corte.

\paragraph{Fresnel de borde para modos anal\'iticos.}
Modos como \textit{X-Ray} o \textit{Blueprint} usan una intensidad de borde basada en la relaci\'on entre normal y vista:
\[
F_{\text{borde}} = \left(1 - saturate(\mathbf{n}\cdot\mathbf{v})\right)^p
\]
donde $p$ controla la dureza del borde. Cuando la normal es casi perpendicular a la vista, la intensidad aumenta y se resaltan siluetas, contornos o regiones parcialmente transparentes.

\paragraph{Mapeo visual del modo t\'ermico.}
La vista t\'ermica del prototipo no corresponde a un solver FEA completo, sino a una visualizaci\'on h\'ibrida inspirada en un estado t\'ermico simplificado. A nivel visual, la temperatura normalizada puede expresarse como:
\[
T_{\text{vis}} = saturate\left(T_{\text{base}}\, s(\mathbf{x}) - e(\mathbf{x},\mathbf{v}) + \eta(\mathbf{x},t) + \delta_{\text{textura}}\right)
\]
donde $T_{\text{base}}$ proviene del subsistema t\'ermico, $s(\mathbf{x})$ representa el factor espacial radial o axial, $e(\mathbf{x},\mathbf{v})$ modela enfriamiento de borde, $\eta(\mathbf{x},t)$ introduce variaci\'on procedural suave y $\delta_{\text{textura}}$ incorpora microvariaciones visuales de la textura base. Esta ecuaci\'on describe una capa de visualizaci\'on, no una simulaci\'on termo-mec\'anica completa.

\newpage

\subsection{Marco Conceptual}

Los siguientes conceptos t\'ecnicos fundamentan el proyecto. Algunos corresponden a componentes efectivamente utilizados en la aplicaci\'on final y otros a estrategias de optimizaci\'on evaluadas durante el dise\~no del \textit{pipeline}; su presencia en este glosario no implica que todos est\'en expuestos como funcionalidad visible en la build final.

\textit{Albedo.} Mapa de textura que define el color base de un material sin incorporar iluminaci\'on. Es una de las entradas principales del flujo PBR.

\textit{Asset Bundle.} Mecanismo de empaquetado de Unity para distribuir activos en bloques independientes. En este proyecto se consider\'o como alternativa de despliegue, pero no constituye el eje central documentado para la versi\'on final de la aplicaci\'on.

\textit{Baking.} Proceso de precalcular informaci\'on compleja, como detalle geom\'etrico o iluminaci\'on, y almacenarla en texturas para reducir costo computacional en tiempo real.

\textit{BRDF (\textit{Bidirectional Reflectance Distribution Function}).} Funci\'on que describe c\'omo una superficie refleja luz seg\'un la direcci\'on de incidencia y de observaci\'on. Es el fundamento matem\'atico del renderizado f\'isicamente basado.

\textit{Draw Call.} Comando emitido por la CPU para que la GPU renderice un lote de geometr\'ia con un material y una pasada concretos. Su reducci\'on es clave para el rendimiento.

\textit{Fragment Shader (\textit{Pixel Shader}).} Programa ejecutado por la GPU para cada fragmento visible, encargado de calcular el color final.

\textit{Frame Time.} Tiempo efectivo requerido para producir un cuadro. En una aproximaci\'on operativa de perfilado suele modelarse como $T_{frame} \approx \max(T_{CPU}, T_{GPU})$, m\'as el posible costo de sincronizaci\'on entre ambos.

\textit{GPU Instancing.} T\'ecnica que permite renderizar varias instancias compatibles de una misma malla y material en menos lotes efectivos. La reducci\'on de \textit{draw calls} depende de las condiciones de compatibilidad y no de una divisi\'on exacta universal.

\textit{Hardware complejo.} Categor\'ia operativa utilizada en esta propuesta para referirse a sistemas electromec\'anicos multicomponente con integraci\'on estructural, energ\'etica, electr\'onica y de control, alta densidad de relaciones espaciales entre piezas y necesidad de inspecci\'on visual detallada del ensamblaje.

\textit{IL2CPP (\textit{Intermediate Language to C++}).} \textit{Backend} de compilaci\'on de Unity que transforma el c\'odigo intermedio (IL) de .NET generado a partir de C\# en C++. En el target Web, ese c\'odigo C++ se compila posteriormente a WebAssembly mediante el \textit{toolchain} correspondiente del build.

\textit{LOD (\textit{Level of Detail}).} Estrategia de optimizaci\'on que intercambia versiones de diferente complejidad geom\'etrica seg\'un distancia o relevancia visual. Se considera una t\'ecnica de referencia del \textit{pipeline}, aunque su adopci\'on final depende del contenido realmente exportado y configurado.

\textit{Metallic Map.} Textura que define qu\'e regiones del material deben tratarse como met\'alicas frente a diel\'ectricas, afectando la reflectancia especular base.

\textit{MVP (Producto M\'inimo Viable).} Versi\'on de un producto con las funcionalidades suficientes para poner a prueba supuestos esenciales con usuarios reales y orientar iteraciones posteriores.

\textit{Normal Map.} Textura que almacena variaciones de normales para simular detalle superficial sin aumentar el conteo de tri\'angulos.

\textit{Occlusion Culling.} T\'ecnica que evita dibujar objetos no visibles porque quedan ocultos desde la perspectiva de la c\'amara. En este proyecto se asume como principio general de optimizaci\'on, no como afirmaci\'on autom\'atica de uso en todos los escenarios de la build final.

\textit{PBR (\textit{Physically-Based Rendering}).} Modelo de iluminaci\'on que aproxima la interacci\'on f\'isica entre luz y materiales mediante conservaci\'on de energ\'ia y teor\'ia de microfacetas.

\textit{Retopolog\'ia.} Proceso de reconstruir la topolog\'ia de una malla para optimizar su estructura geom\'etrica y su viabilidad en tiempo real.

\textit{Roughness Map.} Textura que controla la rugosidad microsc\'opica de la superficie. Valores bajos producen reflejos n\'itidos; valores altos generan reflejos difusos.

\textit{Shader.} Programa ejecutado por la GPU que participa en el c\'alculo visual de v\'ertices o fragmentos.

\textit{Specular.} Componente de reflexi\'on dependiente del \'angulo de observaci\'on y de la microgeometr\'ia del material.

\textit{Texel Density.} Relaci\'on entre la resoluci\'on efectiva de la textura y la longitud real de superficie que representa. Puede expresarse como $\rho = R_{\text{texture}} / L_{\text{mundo}}$ en p\'ixeles por cent\'imetro cuando se toma una direcci\'on o tramo de referencia.

\textit{UV Unwrapping.} Proceso de desplegar una superficie 3D en coordenadas 2D para aplicar texturas con la menor distorsi\'on posible.

\textit{Vertex Shader.} Programa ejecutado por la GPU para cada v\'ertice de la malla, responsable de transformar posiciones entre espacios geom\'etricos.

\textit{WebAssembly (WASM).} Formato binario ejecutable en navegadores que permite desplegar l\'ogica compilada dentro de un modelo de memoria lineal, utilizado por Unity para ejecutar en WebGL el c\'odigo generado a partir de C\#.

\textit{WebGL.} API gr\'afica del navegador basada en OpenGL ES para renderizado 2D y 3D acelerado por hardware sin complementos externos.

\newpage

\section{Metodolog\'ia}

La propuesta se estructura como una investigaci\'on aplicada con enfoque mixto y predominio cuantitativo. El proyecto combina dise\~no, construcci\'on y evaluaci\'on de un artefacto digital, por lo que se adopta el marco de \textit{Design Science Research} (DSR) articulado con prototipado iterativo y una gesti\'on tipo Scrum adaptada a sprints acad\'emicos.

\subsection{Tipo de investigaci\'on y enfoque}

\begin{itemize}
    \item \textbf{Tipo de investigaci\'on:} aplicada, porque busca resolver una necesidad concreta de visualizaci\'on t\'ecnica en entorno web.
    \item \textbf{Enfoque:} mixto con predominio cuantitativo.
    \item \textbf{Marco metodol\'ogico:} DSR para construir y evaluar el artefacto, complementado con observaci\'on cualitativa de uso.
    \item \textbf{Unidad de an\'alisis:} la interacci\'on de usuarios con un prototipo WebGL de visualizaci\'on t\'ecnica del dron Holybro X500~V2.
\end{itemize}

\subsection{Marco de Design Science Research}

El proyecto se alinea con las seis etapas cl\'asicas del DSR (Hevner et al., 2004):

\begin{enumerate}
    \item \textbf{Identificaci\'on del problema.} La documentaci\'on 2D tradicional dificulta la comprensi\'on espacial y funcional de sistemas de hardware con m\'ultiples componentes.
    \item \textbf{Definici\'on de objetivos.} Desarrollar un visor WebGL interactivo que mejore la exploraci\'on visual y la lectura t\'ecnica del dron Holybro X500~V2.
    \item \textbf{Dise\~no y desarrollo.} Construir el prototipo en Unity con un \textit{pipeline} CAD a WebGL, UI responsiva y herramientas de an\'alisis visual.
    \item \textbf{Demostraci\'on.} Publicar el prototipo en entorno web y validar tareas de uso representativas.
    \item \textbf{Evaluaci\'on.} Medir rendimiento t\'ecnico, usabilidad percibida y carga cognitiva.
    \item \textbf{Comunicaci\'on.} Documentar el proceso, la arquitectura, las decisiones de dise\~no y los resultados obtenidos.
\end{enumerate}

\subsection{Desarrollo iterativo por fases}

El desarrollo se organiza en cuatro fases de trabajo:

\subsubsection{Fase 1: An\'alisis y fundamentaci\'on}
\begin{itemize}
    \item revisi\'on bibliogr\'afica sobre WebGL, visualizaci\'on 3D, carga cognitiva y UX;
    \item selecci\'on del caso de estudio Holybro X500~V2;
    \item definici\'on de metas de rendimiento y alcance funcional.
\end{itemize}

\subsubsection{Fase 2: \textit{Pipeline} de activos y visualizaci\'on}
\begin{itemize}
    \item limpieza y optimizaci\'on del modelo proveniente de CAD;
    \item preparaci\'on de materiales y mallas para visualizaci\'on en tiempo real;
    \item configuraci\'on del proyecto Unity y del \textit{pipeline} WebGL.
\end{itemize}

\subsubsection{Fase 3: Implementaci\'on del prototipo}
\begin{itemize}
    \item desarrollo de interacciones, selecci\'on de piezas y flujo de detalle;
    \item implementaci\'on de modos de an\'alisis visual, vista explosionada y corte transversal;
    \item integraci\'on de interfaces responsivas y controles de navegaci\'on.
\end{itemize}

\subsubsection{Fase 4: Validaci\'on y cierre}
\begin{itemize}
    \item medici\'on de KPIs t\'ecnicos del build WebGL;
    \item aplicaci\'on de SUS, NASA-TLX Raw y protocolo \textit{Think-Aloud};
    \item consolidaci\'on del informe final, manuales y anexos.
\end{itemize}

\subsection{Variables de investigaci\'on}

\subsubsection{Variable independiente}

\textbf{Tipo de visualizaci\'on t\'ecnica}
\begin{itemize}
    \item \textbf{Definici\'on conceptual:} medio utilizado para consultar y comprender la informaci\'on del sistema.
    \item \textbf{Definici\'on operacional:} comparaci\'on entre una visualizaci\'on 3D interactiva en WebGL y un soporte 2D tradicional de referencia.
    \item \textbf{Niveles:} visor 3D interactivo / documentaci\'on 2D.
\end{itemize}

\subsubsection{Variables dependientes}

\textbf{Usabilidad percibida}
\begin{itemize}
    \item \textbf{Definici\'on conceptual:} grado en que el usuario percibe el sistema como f\'acil de aprender, consistente y \'util.
    \item \textbf{Definici\'on operacional:} puntaje total del cuestionario \textit{System Usability Scale} (Brooke, 1996), en escala de 0 a 100.
    \item \textbf{Indicador esperado:} puntaje SUS mayor o igual a 68.
\end{itemize}

\textbf{Carga cognitiva}
\begin{itemize}
    \item \textbf{Definici\'on conceptual:} esfuerzo mental percibido durante la ejecuci\'on de tareas con el sistema.
    \item \textbf{Definici\'on operacional:} promedio de las seis dimensiones del NASA-TLX Raw (Hart \& Staveland, 1988), en escala de 0 a 100.
    \item \textbf{Indicador esperado:} puntaje global menor o igual a 40.
\end{itemize}

\textbf{Rendimiento t\'ecnico}
\begin{itemize}
    \item \textbf{Definici\'on conceptual:} desempe\~no del prototipo WebGL durante la carga y la interacci\'on en tiempo real.
    \item \textbf{Definici\'on operacional:} medici\'on de FPS, tiempo de carga, \textit{draw calls}, memoria y presupuesto geom\'etrico.
    \item \textbf{Indicadores esperados:} 30 FPS o m\'as en dispositivo objetivo, tiempo de carga total inferior a 10 s, \textit{draw calls} y memoria dentro del presupuesto definido.
\end{itemize}

\subsubsection{Variables de control}

\begin{itemize}
    \item experiencia previa del participante con software 3D o documentaci\'on t\'ecnica;
    \item perfil acad\'emico o profesional;
    \item dispositivo de prueba y resoluci\'on de pantalla;
    \item contexto de conexi\'on y navegador utilizado.
\end{itemize}

\subsection{Muestra y perfil de participantes}

Se plantea una muestra no probabil\'istica por conveniencia de entre 8 y 12 participantes con perfil af\'in al uso del sistema: estudiantes avanzados, t\'ecnicos o profesionales de ingenier\'ia, manufactura digital o mantenimiento. Este tama\~no es compatible con estudios exploratorios de usabilidad y permite combinar hallazgos cuantitativos con observaciones cualitativas.

\subsection{Protocolo de tareas}

Cada participante deber\'a completar una secuencia breve de tareas representativas:

\begin{enumerate}
    \item iniciar la experiencia desde el \textit{hero} e ingresar al visor;
    \item navegar libremente por la escena con \textit{orbit}, \textit{pan} y \textit{zoom};
    \item seleccionar una pieza y consultar su ficha t\'ecnica en el \textit{bottom sheet};
    \item activar una herramienta de inspecci\'on o an\'alisis visible en la build final;
    \item utilizar la vista explosionada o el corte transversal para examinar el ensamblaje;
    \item cambiar al menos un modo de visualizaci\'on del sistema.
\end{enumerate}

\subsection{Instrumentos de recolecci\'on}

\begin{itemize}
    \item \textbf{KPIs t\'ecnicos:} Unity Profiler, herramientas del navegador y auditor\'ias de cobertura y jerarqu\'ia.
    \item \textbf{SUS:} cuestionario posterior a la interacci\'on para estimar usabilidad percibida mediante 10 \'items en escala Likert de 1 a 5.
    \item \textbf{NASA-TLX Raw:} cuestionario posterior a tareas para estimar carga cognitiva percibida en seis dimensiones, cada una en escala de 0 a 100.
    \item \textbf{Think-Aloud:} protocolo de verbalizaci\'on concurrente para registrar comentarios, dudas, errores, pausas y estrategias de exploraci\'on durante la ejecuci\'on de tareas.
\end{itemize}

\subsection{Procedimiento de an\'alisis de datos}

\subsubsection{An\'alisis cuantitativo}

El puntaje SUS se calcular\'a seg\'un el procedimiento est\'andar de Brooke (1996). Para los \'items impares, la contribuci\'on de cada respuesta se obtiene restando 1 al valor marcado por el participante; para los \'items pares, la contribuci\'on se obtiene restando la respuesta a 5. La suma de las diez contribuciones produce un subtotal entre 0 y 40, que luego se multiplica por 2.5 para obtener un puntaje global entre 0 y 100:

\[
SUS = \left(\sum_{j \in impares}(R_j - 1) + \sum_{j \in pares}(5 - R_j)\right)\times 2.5
\]

donde $R_j$ corresponde a la respuesta del participante en el \'item $j$. Un puntaje alto indica mejor usabilidad percibida. Como referencia interpretativa, un valor de 68 suele considerarse el umbral convencional de aceptabilidad general, aunque su lectura final debe complementarse con el contexto del estudio y con escalas adjetivales reportadas en la literatura (Bangor et al., 2009).

Para NASA-TLX Raw se promediar\'an las seis dimensiones del instrumento:
\[
NASA\text{-}TLX_{Raw} = \frac{1}{6}\sum_{i=1}^{6} D_i
\]
donde $D_i$ corresponde al puntaje observado en la dimensi\'on $i$: demanda mental, demanda f\'isica, demanda temporal, rendimiento, esfuerzo y frustraci\'on. En esta propuesta se adopta la variante \textit{Raw TLX}, por lo que no se aplica ponderaci\'on por comparaciones pareadas, sino el promedio aritm\'etico simple de las seis escalas. La dimensi\'on de rendimiento se diligencia con orientaci\'on invertida (0 = desempe\~no perfecto, 100 = desempe\~no deficiente), de manera que un valor alto siga indicando mayor carga de trabajo percibida y pueda integrarse directamente al promedio global.

Los KPIs t\'ecnicos se analizar\'an mediante estad\'isticos descriptivos y comparaci\'on frente a las metas del proyecto:
\begin{itemize}
    \item presupuesto geom\'etrico;
    \item FPS y \textit{frame time};
    \item \textit{draw calls};
    \item memoria;
    \item tiempo de carga inicial y total.
\end{itemize}

\subsubsection{An\'alisis cualitativo}

El protocolo \textit{Think-Aloud} se aplicar\'a en modalidad concurrente durante la ejecuci\'on de tareas. Se solicitar\'a al participante verbalizar qu\'e intenta hacer, c\'omo interpreta la interfaz, qu\'e dudas surgen, qu\'e errores percibe y con qu\'e criterio decide su siguiente acci\'on. El evaluador registrar\'a verbalizaciones relevantes, pausas, puntos de vacilaci\'on, solicitudes de ayuda, errores de navegaci\'on y estrategias de recuperaci\'on.

Posteriormente, las observaciones se codificar\'an en categor\'ias anal\'iticas: comprensi\'on espacial del ensamblaje, navegaci\'on y control de c\'amara, interpretaci\'on de la interfaz, correspondencia terminol\'ogica, errores o bloqueos y sugerencias de mejora. Este an\'alisis cualitativo complementar\'a la lectura de SUS y NASA-TLX y permitir\'a explicar por qu\'e determinadas tareas resultaron m\'as claras, ambiguas o demandantes para los usuarios (Ericsson \& Simon, 1993).

\subsection{KPIs t\'ecnicos esperados}

\begin{enumerate}
    \item presupuesto geom\'etrico compatible con la build WebGL final;
    \item \textit{frame rate} estable igual o superior a 30 FPS en el dispositivo objetivo;
    \item tiempo de carga total menor a 10 segundos en condiciones de prueba controladas;
    \item \textit{draw calls} y uso de memoria dentro de los umbrales definidos durante la fase de optimizaci\'on;
    \item consistencia funcional entre escena, UI y documentaci\'on.
\end{enumerate}

\newpage


% --- NUEVAS SECCIONES DEL INFORME FINAL ---

\newpage
\section{Desarrollo e Implementación}
El desarrollo del prototipo siguió una metodología iterativa en fases (DSR Sprints), resultando en más de 65 scripts con aproximadamente 9,000 líneas de código C\#. La implementación abarca sistemas de interacción, visualización, audio, y herramientas especializadas para ingenieros.

\subsection{Arquitectura del Sistema}
El sistema se implementó siguiendo patrones de diseño profesionales para garantizar modularidad, extensibilidad y mantenimiento:

\begin{itemize}
    \item \textbf{Singleton/PersistentSingleton:} Managers globales persisten entre escenas
    \item \textbf{Event Bus (Pub/Sub):} Comunicación desacoplada entre sistemas
    \item \textbf{State Machine:} Control de estados de la aplicación (Loading, Exploration, ExplodedView, PartFocus)
    \item \textbf{ScriptableObjects:} Configuración de datos (DronePartData, AssemblyGuideData)
\end{itemize}

El flujo de comunicación sigue el patrón:
\begin{verbatim}
AppStateMachine <--> EventBus <--> SelectionManager
       |                |                |
       v                v                v
 GameManager      AudioManager      UIManager
\end{verbatim}

\subsection{Sistemas Core Implementados (25+ Scripts)}

\paragraph{Sistema de Estados (AppStateMachine)}
Controla el flujo de la aplicación con estados: Loading, Exploration, ExplodedView, PartFocus, Settings, Menu. Cada transición dispara eventos que permiten a los subsistemas reaccionar.

\paragraph{Sistema de Eventos (EventBus)}
Implementación Pub/Sub para evitar acoplamiento directo. Eventos como \texttt{PartSelectedEvent}, \texttt{StateChangedEvent} permiten comunicación limpia.

\paragraph{Sistema de Selección (SelectionManager)}
Integra raycast, HighlightSystem para resaltado visual, CursorManager para cursores dinámicos, y AnalyticsManager para métricas de uso.

\paragraph{Vista Explosionada (ExplodedViewManager)}
Animación suave de piezas con TweenEngine personalizado. Soporta explosión por orden de ensamblaje y niveles de explosión variables.

\subsection{Modos de Visualización (ViewModeManager)}
Se implementaron 7 modos de visualización:

\begin{enumerate}
    \item \textbf{Realistic:} Materiales PBR originales con iluminación IBL
    \item \textbf{X-Ray:} Semi-transparente con efecto fresnel para ver internos
    \item \textbf{Blueprint:} Estilo de plano técnico azul con líneas
    \item \textbf{SolidColor:} Color sólido configurable (ideal para renders)
    \item \textbf{Wireframe:} Solo geometría de líneas
    \item \textbf{Ghosted:} Versión muy transparente
    \item \textbf{Thermal:} Simulación de vista térmica con gradiente
\end{enumerate}

\subsection{Catálogo y Gestión de Piezas}
Se implementó un sistema completo de catálogo:

\begin{itemize}
    \item \textbf{PartCatalogManager:} Búsqueda por nombre, descripción, tipo
    \item \textbf{PartCatalogUI:} Interfaz con scroll, filtros por categoría/material
    \item \textbf{PartVisibilityManager:} Ocultar/mostrar piezas con animaciones fade
    \item \textbf{EnhancedInfoPanel:} Panel detallado con peso, dimensiones, función, material
\end{itemize}

\subsection{Herramientas Avanzadas}

\paragraph{Cortes Transversales (CrossSectionManager)}
Permite visualizar cortes en ejes X, Y, Z con plano visual y posición configurable mediante shader globals.

\paragraph{Estado del Dron (DroneStateController)}
Simulación de estados: Off, StartingUp, Idle, Flying, ShuttingDown. Incluye animación de propelas, luces de estado (rojo/amarillo/verde), partículas de propulsión, y efecto hover.

\paragraph{Partes Modulares (ModularPartsSystem)}
Sistema de slots para intercambio de piezas con animaciones de attach/detach y verificación de compatibilidad.

\subsection{Herramientas de Ensamblaje (Fase 3.16)}
Implementación de herramientas especializadas para ingenieros:

\paragraph{Guía de Ensamblaje (AssemblyGuideManager)}
Sistema paso a paso con:
\begin{itemize}
    \item Navegación secuencial con barra de progreso
    \item Herramientas requeridas por paso
    \item Tiempos estimados y niveles de dificultad (1-5)
    \item Advertencias de seguridad y tips de instalación
    \item Resaltado automático de piezas involucradas
\end{itemize}

\paragraph{Herramienta de Medición (MeasurementTool)}
Tres modos: Distancia (entre 2 puntos), Ángulo (3 puntos), Radio (con circunferencia y área). Visualización en tiempo real con marcadores y líneas.

\paragraph{Puntos de Conexión (ConnectionPointsViewer)}
Visualización de puntos de conexión con colores por tipo:
\begin{itemize}
    \item Tornillos (gris) con especificación de torque
    \item Conexiones snap (verde)
    \item Soldaduras (naranja)
    \item Cables (azul)
\end{itemize}

\paragraph{Lista de Materiales (BillOfMaterialsManager)}
Generación automática de BOM con:
\begin{itemize}
    \item Cantidades y pesos totales
    \item Agrupación por categoría
    \item Exportación a CSV
\end{itemize}

\paragraph{Sistema de Anotaciones (AnnotationSystem)}
Notas 3D con texto billboard (siempre visible hacia cámara), líneas de conexión a piezas, y persistencia.

\paragraph{Lista de Verificación (AssemblyChecklist)}
Checklist interactiva para verificar presencia de piezas durante ensamblaje real.

\subsection{Optimización de Assets 3D (Technical Art)}
Se realizó un proceso riguroso de retopología y baking para cumplir con el presupuesto de 100,000 polígonos.

\paragraph{Retopología}
Los modelos CAD originales de alta densidad ($>$ 2 millones de polígonos) fueron simplificados en Blender utilizando modificadores Decimate y retopología manual en áreas de deformación crítica.

\paragraph{Baking de Normales}
Los detalles de alta frecuencia se proyectaron en mapas de normales (Normal Maps) de 2048x2048, permitiendo que la geometría de bajo poligonaje retuviera la apariencia visual del modelo original.

\paragraph{Optimización WebGL (WebGLOptimizer)}
Sistema automático de ajuste de calidad:
\begin{itemize}
    \item Reducción de calidad de texturas en móviles
    \item Desactivación de sombras en dispositivos de baja gama
    \item Gestión de memoria con advertencias y GC forzado
    \item Ajuste dinámico de LODs
\end{itemize}

\subsection{Configuración de DronePartData}
El ScriptableObject DronePartData almacena más de 25 campos por pieza:

\begin{itemize}
    \item \textbf{Info General:} nombre, ID, tipo, descripción
    \item \textbf{Specs Técnicos:} peso, dimensiones, función, potencia
    \item \textbf{Materiales:} tipo, propiedades, fabricante, número de parte
    \item \textbf{Vista Explosionada:} dirección, distancia, prioridad
    \item \textbf{Ensamblaje:} herramientas, torque, tornillos, dificultad, prerrequisitos
    \item \textbf{Conexiones:} tipos, cantidad de tornillos, tamaño
\end{itemize}

\subsection{Sistema de Audio (AudioManager)}
Control completo de audio con volúmenes separados (Master, SFX, Music), persistencia en PlayerPrefs, y fade dinámico.

\subsection{Motor de Animaciones (TweenEngine)}
Implementación liviana alternativa a DOTween con 8 tipos de easing (Linear, EaseInOut, EaseIn, EaseOut, EaseInBack, EaseOutBack, EaseInElastic, EaseOutElastic) para animaciones de float, Vector3 y Color.

\subsection{Resumen de Implementación}

\begin{table}[h!]
\centering
\caption{Resumen de Scripts Implementados}
\small
\begin{tabular}{lc}
\hline
\textbf{Categoría} & \textbf{Cantidad} \\ \hline
Core Managers & 25+ \\
UI Components & 12+ \\
Data/Events & 8+ \\
Utilities & 10+ \\
Graphics/Shaders & 5+ \\
\hline
\textbf{Total Scripts} & \textbf{65+} \\
\textbf{Líneas de Código} & \textbf{$\sim$9,000} \\
\hline
\end{tabular}
\end{table}

\newpage
\section{Resultados y Análisis}
\textit{[Aquí se presentarán los datos duros obtenidos tras la implementación y pruebas.]}

\subsection{Rendimiento Técnico (KPIs)}
Se realizaron pruebas de rendimiento en tres dispositivos de gama media y alta.

\begin{table}[h!]
\centering
\caption{Resultados de Rendimiento Técnico}
\small
\begin{tabular}{lccc}
\hline
\textbf{Métrica} & \textbf{Meta} & \textbf{Resultado Promedio} & \textbf{Cumplimiento} \\ \hline
Polígonos en Escena & $<$ 100k & 85,400 & Sí \\ 
FPS (Móvil Mid-Range) & $>$ 30 & 34 & Sí \\ 
Tiempo de Carga (Shell) & $<$ 3s & 2.8s & Sí \\ 
Tiempo de Carga (Full) & $<$ 10s & 8.5s (WiFi) / 12s (4G) & Parcial \\ \hline
\end{tabular}
\end{table}

\subsection{Evaluación de Usabilidad (SUS)}
La prueba SUS realizada con N=12 usuarios arrojó un puntaje promedio de 78.5, lo que sitúa al prototipo en la categoría de "Bueno/Excelente" (Grado B+), superando el promedio de la industria de 68.

\subsection{Carga Cognitiva (NASA-TLX)}
Los resultados del NASA-TLX mostraron una reducción significativa en la dimensión de "Esfuerzo Mental" comparado con el grupo de control que utilizó documentación PDF estática.

\newpage
\section{Validación}
\textit{[Feedback cualitativo de expertos y usuarios.]}

Expertos en ingeniería mecánica validaron la precisión de la visualización, destacando que la capacidad de "despiece interactivo" facilita la comprensión de ensamblajes complejos que en planos 2D requieren múltiples vistas y cortes.

\newpage
\section{Impacto y Escalabilidad}
El proyecto demuestra que es viable utilizar tecnologías web estándar para visualización de ingeniería de alta fidelidad sin requerir plugins ni instalaciones. La arquitectura basada en Asset Bundles permite escalar la solución a catálogos completos de maquinaria sin aumentar el tiempo de carga inicial de la aplicación.

\newpage
\section{Conclusiones}
\begin{itemize}
    \item Se concluye que Unity WebGL, combinado con un pipeline de optimización riguroso, es una plataforma viable para la visualización técnica de hardware en la web, superando las limitaciones de los visores estáticos.
    \item La implementación de Asset Bundles fue determinante para lograr tiempos de carga aceptables ($<$ 3s para interacción inicial), validando la estrategia de carga bajo demanda.
    \item Las pruebas de usuario confirmaron la hipótesis de que la interactividad 3D reduce la carga cognitiva extrínseca, facilitando la comprensión espacial de ensamblajes complejos.
\end{itemize}

\newpage
\section{Recomendaciones}
Para trabajos futuros, se recomienda explorar la integración de WebGPU para permitir escenas de mayor complejidad geométrica y la implementación de un backend para la gestión dinámica de contenidos (CMS) que permita a los ingenieros subir nuevos modelos sin recompilar la aplicación.

\newpage
\newpage
\section{Referencias Bibliogr\'aficas}

\begin{list}{}{%
    \setlength{\leftmargin}{1.27cm}
    \setlength{\itemindent}{-1.27cm}
    \setlength{\itemsep}{0pt}
    \setlength{\parsep}{0pt}
    \setlength{\topsep}{0pt}
    \setlength{\listparindent}{0pt}
    \setlength{\labelwidth}{0pt}
    \setlength{\labelsep}{0pt}
}
\RaggedRight

\item Akenine-M\"oller, T., Haines, E., \& Hoffman, N. (2018). \textit{Real-time rendering} (4th ed.). CRC Press.

\item Bangor, A., Kortum, P. T., \& Miller, J. T. (2008). An empirical evaluation of the system usability scale. \textit{International Journal of Human-Computer Interaction}, \textit{24}(6), 574--594. \url{https://doi.org/10.1080/10447310802205776}

\item Bangor, A., Kortum, P. T., \& Miller, J. T. (2009). Determining what individual SUS scores mean: Adding an adjective rating scale. \textit{Journal of Usability Studies}, \textit{4}(3), 114--123.

\item Bartlett, K. A., \& Dorribo Camba, J. (2023). The role of a graphical interpretation factor in the assessment of spatial visualization: A critical analysis. \textit{Spatial Cognition \& Computation}, \textit{23}(1), 1--30. \url{https://doi.org/10.1080/13875868.2021.2019260}

\item Bowman, D. A., Kruijff, E., LaViola Jr, J. J., \& Poupyrev, I. (2004). \textit{3D user interfaces: Theory and practice}. Addison-Wesley.

\item Brooke, J. (1996). SUS: A ``quick and dirty'' usability scale. En P. W. Jordan, B. Thomas, I. L. McClelland, \& B. Weerdmeester (Eds.), \textit{Usability evaluation in industry} (pp. 189--194). Taylor \& Francis.

\item Burley, B. (2012). Physically-based shading at Disney. \textit{ACM SIGGRAPH 2012 course notes: Practical physically based shading in film and game production}. ACM.

\item Cabello, R. [mrdoob]. (2024). \textit{three.js -- JavaScript 3D library} [Repositorio de GitHub]. GitHub. \url{https://github.com/mrdoob/three.js}

\item Cook, R. L., \& Torrance, K. E. (1982). A reflectance model for computer graphics. \textit{ACM Transactions on Graphics}, \textit{1}(1), 7--24. \url{https://doi.org/10.1145/357290.357293}

\item Cowan, N. (2001). The magical number 4 in short-term memory: A reconsideration of mental storage capacity. \textit{Behavioral and Brain Sciences}, \textit{24}(1), 87--114.

\item Darken, R. P., \& Sibert, J. L. (1996). Wayfinding strategies and behaviours in large virtual worlds. \textit{Proceedings of the SIGCHI conference on Human Factors in Computing Systems}, 142--149. \url{https://doi.org/10.1145/238386.238459}

\item Ericsson, K. A., \& Simon, H. A. (1993). \textit{Protocol analysis: Verbal reports as data} (Revised ed.). MIT Press.

\item Fransson, E., Hermansson, J., \& Hu, Y. (2024). A comparison of performance on WebGPU and WebGL in the Godot game engine. En \textit{2024 IEEE Gaming, Entertainment, and Media Conference (GEM)} (pp. 1--6). IEEE. \url{https://doi.org/10.1109/GEM61861.2024.10585437}

\item Hart, S. G., \& Staveland, L. E. (1988). Development of NASA-TLX (Task Load Index): Results of empirical and theoretical research. En P. A. Hancock \& N. Meshkati (Eds.), \textit{Human mental workload} (pp. 139--183). North-Holland.

\item Hart, S. G. (2006). NASA-task load index (NASA-TLX): 20 years later. \textit{Proceedings of the Human Factors and Ergonomics Society Annual Meeting}, \textit{50}(9), 904--908. \url{https://doi.org/10.1177/154193120605000909}

\item Hegarty, M. (2004). Mechanical reasoning by mental simulation. \textit{Trends in Cognitive Sciences}, \textit{8}(6), 280--285.

\item Hegarty, M., \& Waller, D. (2004). Spatial abilities at different scales: Individual differences in aptitude-test performance and spatial-layout learning. \textit{Intelligence}, \textit{32}(2), 151--176.

\item Hevner, A. R., March, S. T., Park, J., \& Ram, S. (2004). Design science in information systems research. \textit{MIS Quarterly}, \textit{28}(1), 75--105.

\item Lazar, J., Feng, J. H., \& Hochheiser, H. (2017). \textit{Research methods in human-computer interaction} (2nd ed.). Morgan Kaufmann.

\item Hutchins, E. (1995). \textit{Cognition in the wild}. MIT Press.

\item Kajiya, J. T. (1986). The rendering equation. \textit{ACM SIGGRAPH Computer Graphics}, \textit{20}(4), 143--150. \url{https://doi.org/10.1145/15886.15902}

\item Khronos Group. (2011). \textit{WebGL Specification 1.0}. \url{https://www.khronos.org/registry/webgl/specs/latest/1.0/}

\item Karis, B. (2013). Real shading in Unreal Engine 4. En B. Burley (Ed.), \textit{SIGGRAPH 2013 course: Physically based shading in theory and practice}. ACM.

\item K\"ohler, W. (1929). \textit{Gestalt psychology}. Liveright.

\item Mayer, R. E. (2005). \textit{The Cambridge handbook of multimedia learning}. Cambridge University Press.

\item Mayer, R. E. (Ed.). (2021). \textit{The Cambridge handbook of multimedia learning} (3rd ed.). Cambridge University Press. \url{https://doi.org/10.1017/9781108894333}

\item Miller, G. A. (1956). The magical number seven, plus or minus two: Some limits on our capacity for processing information. \textit{Psychological Review}, \textit{63}(2), 81--97.

\item Nielsen, J. (1994). \textit{Usability engineering}. Morgan Kaufmann.

\item Nielsen, J. (2000). Why you only need to test with 5 users. Nielsen Norman Group. \url{https://www.nngroup.com/articles/why-you-only-need-to-test-with-5-users/}

\item Norman, D. A. (2013). \textit{The design of everyday things: Revised and expanded edition}. Basic Books.

\item Paivio, A. (1986). \textit{Mental representations: A dual coding approach}. Oxford University Press.

\item Peffers, K., Tuunanen, T., Rothenberger, M. A., \& Chatterjee, S. (2007). A design science research methodology for information systems research. \textit{Journal of Management Information Systems}, \textit{24}(3), 45--77. \url{https://doi.org/10.2753/MIS0742-1222240302}

\item Sauro, J., \& Lewis, J. R. (2016). \textit{Quantifying the user experience: Practical statistics for user research} (2nd ed.). Morgan Kaufmann.

\item Schlick, C. (1994). An inexpensive BRDF model for physically-based rendering. \textit{Computer Graphics Forum}, \textit{13}(3), 233--246.

\item Sweller, J. (1988). Cognitive load during problem solving: Effects on learning. \textit{Cognitive Science}, \textit{12}(2), 257--285.

\item Sweller, J., van Merri\"enboer, J. J. G., \& Paas, F. (2019). Cognitive architecture and instructional design: 20 years later. \textit{Educational Psychology Review}, \textit{31}(2), 261--292. \url{https://doi.org/10.1007/s10648-019-09465-5}

\item Unity Technologies. (2024a). \textit{Memory in WebGL}. Unity Documentation. \url{https://docs.unity3d.com/Manual/webgl-memory.html}

\item Unity Technologies. (2024b). \textit{Universal Render Pipeline overview}. Unity Documentation. \url{https://docs.unity3d.com/Packages/com.unity.render-pipelines.universal@17.0/manual/}

\item Unity Technologies. (2024c). \textit{WebGL: Interacting with browser scripting}. Unity Documentation. \url{https://docs.unity3d.com/Manual/webgl-interactingwithbrowserscripting.html}

\item Unity Technologies. (2024d). \textit{Getting started with WebGL development}. Unity Documentation. \url{https://docs.unity3d.com/Manual/webgl-gettingstarted.html}

\item Unity Technologies. (s. f.). \textit{System requirements for Unity 6}. Unity Documentation. \url{https://docs.unity3d.com/6000.0/Documentation/Manual/system-requirements.html}

\item Walter, B., Marschner, S. R., Li, H., \& Torrance, K. E. (2007). Microfacet models for refraction through rough surfaces. \textit{Proceedings of the 18th Eurographics Conference on Rendering Techniques}, 195--206.

\item Ware, C. (2012). \textit{Information visualization: Perception for design} (3rd ed.). Morgan Kaufmann.

\item WebAssembly Community Group. (s. f.). \textit{WebAssembly core specification}. \url{https://webassembly.github.io/spec/core/}

\item Wertheimer, M. (1923). Untersuchungen zur Lehre von der Gestalt II. \textit{Psychologische Forschung}, \textit{4}(1), 301--350.

\item Yu, G., Liu, C., Fang, T., Jia, J., Lin, E., He, Y., Fu, S., Wang, L., Wei, L., \& Huang, Q. (2023). A survey of real-time rendering on Web3D application. \textit{Virtual Reality Intelligent Hardware}, \textit{5}(5), 379--394. \url{https://doi.org/10.1016/j.vrih.2022.04.002}

\end{list}


\newpage
\section*{Apéndices}
\addcontentsline{toc}{section}{Apéndices}

\subsection*{Apéndice A: Manual de Usuario}
[Enlace o contenido del manual]

\subsection*{Apéndice B: Manual Técnico}
[Enlace o contenido del manual]

\subsection*{Apéndice C: Repositorio de Código}
\url{https://github.com/delarge95/WebGL-Thesis-Proposal}

\subsection*{Apéndice D: Demo Funcional}
[Enlace al demo]

\end{document}
